% !TeX program = xelatex
\documentclass[aspectratio=169,10pt]{beamer}
\usepackage{amsmath,amssymb,array,booktabs,mathtools}
\usepackage{graphicx,hyperref}
\usepackage{tikz}
% ============================================================
% Beamer Color System — Speculative Decoding black theme
% ============================================================

\providecommand{\beamerthemename}{black}

\RequirePackage{xifthen}

% Keep the public theme selector compatible with the source deck. Its current
% variants intentionally resolve to the same monochrome academic palette.
\newcommand{\usemonochrometheme}{%
  \definecolor{themePrimary}{HTML}{111111}%
  \definecolor{themeSecondary}{HTML}{444444}%
  \definecolor{themeTertiary}{HTML}{000000}%
  \definecolor{themeLight}{HTML}{B8B8B8}%
  \definecolor{themeSilver}{HTML}{E6E6E6}%
  \definecolor{themeBg}{HTML}{FAFAFA}%
  \definecolor{themeBorder}{HTML}{D7D7D7}%
  \definecolor{themeAlert}{HTML}{111111}%
  \definecolor{themeExample}{HTML}{111111}%
}

\ifthenelse{\equal{\beamerthemename}{black}}{\usemonochrometheme}{%
\ifthenelse{\equal{\beamerthemename}{blue}}{\usemonochrometheme}{%
\ifthenelse{\equal{\beamerthemename}{gold}}{\usemonochrometheme}{%
\ifthenelse{\equal{\beamerthemename}{green}}{\usemonochrometheme}{%
\ifthenelse{\equal{\beamerthemename}{purple}}{\usemonochrometheme}{%
  \PackageError{beamer-colors}{Unknown theme '\beamerthemename'.%
    Valid themes: black, blue, gold, green, purple}{}
}}}}}

\setbeamercolor{structure}{fg=themePrimary}
\setbeamercolor{palette primary}{bg=themePrimary,fg=white}
\setbeamercolor{palette secondary}{bg=themeSecondary,fg=white}
\setbeamercolor{palette tertiary}{bg=themeTertiary,fg=white}
\setbeamercolor{palette quaternary}{bg=themePrimary!80!black,fg=white}

\setbeamercolor{title}{fg=themePrimary}
\setbeamercolor{frametitle}{fg=themePrimary}
\setbeamercolor{subtitle}{fg=themeSecondary}
\setbeamercolor{author}{fg=themePrimary!80!black}
\setbeamercolor{institute}{fg=themeSecondary}
\setbeamercolor{date}{fg=themeSecondary}

\setbeamercolor{normal text}{fg=black!84,bg=white}
\setbeamercolor{alerted text}{fg=themeAlert}
\setbeamercolor{example text}{fg=themeExample}

\setbeamercolor{block title}{bg=themePrimary,fg=white}
\setbeamercolor{block body}{bg=themeSilver!40}
\setbeamercolor{block title alerted}{bg=themeAlert,fg=white}
\setbeamercolor{block body alerted}{bg=themeAlert!8}
\setbeamercolor{block title example}{bg=themeExample,fg=white}
\setbeamercolor{block body example}{bg=themeExample!8}

\setbeamercolor{itemize item}{fg=themePrimary}
\setbeamercolor{itemize subitem}{fg=themeSecondary}
\setbeamercolor{itemize subsubitem}{fg=themeSecondary!70}
\setbeamercolor{enumerate item}{fg=themePrimary}
\setbeamercolor{enumerate subitem}{fg=themeSecondary}

\setbeamercolor{section in toc}{fg=themePrimary}
\setbeamercolor{subsection in toc}{fg=themeSecondary}
\setbeamercolor{footnote mark}{fg=themePrimary}
\setbeamercolor{caption name}{fg=themePrimary}

\definecolor{positive}{HTML}{111111}
\definecolor{negative}{gray}{0.42}
\definecolor{emphasis}{HTML}{111111}
\definecolor{neutral}{gray}{0.55}

\newcommand{\pos}[1]{\textcolor{positive}{#1}}
\newcommand{\con}[1]{\textcolor{negative}{#1}}
\newcommand{\HL}[1]{\textcolor{emphasis}{#1}}

\makeatletter
\@ifpackageloaded{hyperref}{%
  \hypersetup{
    colorlinks=true,
    linkcolor=black!75,
    urlcolor=black!75,
    citecolor=black!75,
  }%
}{}
\makeatother

% ============================================================
% Layout: Metropolis — Modern minimal, no external dependencies
% Recreates the metropolis aesthetic using standard beamer.
% ============================================================

\usetheme{default}
\usecolortheme{default}

\newlength{\bsProgressDone}
\newlength{\bsProgressRemain}

% ---- Typography ----
\usefonttheme{professionalfonts}
\setbeamerfont{title}{size=\Large,series=\bfseries}
\setbeamerfont{frametitle}{size=\large,series=\bfseries}
\setbeamerfont{framesubtitle}{size=\normalsize,series=\mdseries}
\setbeamerfont{author}{size=\normalsize}
\setbeamerfont{institute}{size=\small}
\setbeamerfont{date}{size=\small}
\setbeamerfont{section title}{size=\Large,series=\bfseries}
\setbeamerfont{footline}{size=\tiny}

% ---- Remove navigation ----
\setbeamertemplate{navigation symbols}{}

% ---- Frame title: left-aligned, clean ----
\setbeamertemplate{frametitle}{%
  \vskip6pt%
  \usebeamerfont{frametitle}\usebeamercolor[fg]{frametitle}\insertframetitle%
  \ifx\insertframesubtitle\empty\else
    \\[2pt]\usebeamerfont{framesubtitle}\usebeamercolor[fg]{subtitle}\insertframesubtitle%
  \fi
  \vskip4pt%
  {\color{themePrimary}\hrule height 0.45pt}%
  \vskip6pt%
}

% ---- Footline: thin progress bar + page number ----
\setbeamertemplate{footline}{%
  \nointerlineskip%
  \pgfmathsetmacro{\progressfrac}{\insertframenumber/\inserttotalframenumber}%
  \pgfmathsetlength{\bsProgressDone}{\progressfrac\paperwidth}%
  \pgfmathsetlength{\bsProgressRemain}{\paperwidth-\bsProgressDone}%
  \hbox to\paperwidth{%
    {\color{themePrimary!75}\rule{\bsProgressDone}{1pt}}%
    {\color{themeBorder!65}\rule{\bsProgressRemain}{1pt}}%
  }%
  \vskip2pt%
  \hbox to\paperwidth{\hfill\usebeamerfont{footline}\textcolor{black!45}{\insertframenumber\,/\,\inserttotalframenumber}\hspace{8pt}}%
  \vskip4pt%
}

% ---- Title page: minimal ----
\setbeamertemplate{title page}{%
  \vfill
  \begin{flushleft}
    {\usebeamerfont{title}\usebeamercolor[fg]{title}\inserttitle}\\[18pt]
    {\usebeamerfont{author}\textcolor{black!60}{\insertauthor}}\\[4pt]
    {\usebeamerfont{date}\textcolor{black!60}{\insertdate}}
  \end{flushleft}
  \vfill
}

% ---- Section page ----
\AtBeginSection{%
  \begin{frame}[plain,noframenumbering]
    \vfill
    \begin{center}
      {\usebeamerfont{section title}\usebeamercolor[fg]{structure}\insertsectionhead}
    \end{center}
    \vfill
  \end{frame}
}

% ---- Itemize: clean circles ----
\setbeamertemplate{itemize item}{\small\raise1pt\hbox{\textbullet}}
\setbeamertemplate{itemize subitem}{\scriptsize\raise1pt\hbox{--}}
\setbeamertemplate{itemize subsubitem}{\scriptsize\raise1pt\hbox{$\cdot$}}

\author{Skill QA}
\date{5 September 2026}
\newcommand{\slidecite}[1]{%
  \begin{tikzpicture}[remember picture,overlay]
    \node[anchor=south west,align=left,text width=0.92\paperwidth,
      inner sep=0pt,font=\tiny,text=black!38]
      at ([xshift=0.68cm,yshift=0.72cm]current page.south west)
      {\raggedright #1\par};
  \end{tikzpicture}%
}
\newcommand{\norm}[1]{\left\lVert #1\right\rVert}

\title{Adaptive Draft Length Under Batch Load:\\A Proposal}
\hypersetup{pdftitle={Adaptive Draft Length Under Batch Load: A Proposal},pdfauthor={Skill QA}}
\begin{document}
\begin{frame}[plain]
  \titlepage
\end{frame}

\begin{frame}{Batching may change verification cost}
  \vspace{0.2cm}
  \begin{center}
    {\large
      \(\text{draft }K\text{ candidates}
      \ \longrightarrow\ \text{verify in batch }B
      \ \longrightarrow\ \text{emit tokens}\)}
  \end{center}
  \vspace{0.35cm}
  \(K\): proposed tokens per request.\qquad \(B\): active requests in the batch.
  \vspace{0.4cm}
  \begin{itemize}
    \item More candidates may change both verification work and accepted output.
    \item Request latency and total throughput may favor different choices.
  \end{itemize}
  \vspace{0.3cm}
  {\color{black!55}Supplied concern only. No cost curve or measured speedup yet.}
  \slidecite{Supplied proposal fixture P1; toy process M1 in source-notes.md. These are questions to test.}
\end{frame}

\begin{frame}{Every comparison must preserve the target distribution}
  \vspace{0.15cm}
  \[
    \underbrace{q(\,\cdot\mid c)}_{\text{output with the length policy}}
    =
    \underbrace{p(\,\cdot\mid c)}_{\text{target-only output}}
    \qquad\text{for each prompt }c.
  \]
  \vspace{0.3cm}
  \begin{itemize}
    \item Keep models, verifier and decoding settings fixed across policies.
    \item Check full distributions in an enumerable toy, then run sampled regressions.
  \end{itemize}
  \vspace{0.35cm}
  \begin{center}
    A fixed draft length is the baseline.\\[0.15cm]
    Finite checks do not prove correctness for every prompt.
  \end{center}
  \slidecite{Required correctness checks from P1; proposed evaluation contract C1 in source-notes.md.}
\end{frame}

\begin{frame}{The useful draft length may depend on load}
  \vspace{0.1cm}
  {\large Hypothesis}\qquad The lowest cost per emitted token may occur at different \(K\) for different \(B\).
  \vspace{0.45cm}
  \begin{center}
    \renewcommand{\arraystretch}{1.5}
    \begin{tabular}{lp{8.2cm}}
      \toprule
      Proposed step & Length policy\\
      \midrule
      Calibrate & Estimate time per emitted token for each \((B,K)\).\\
      Select & Observe \(B\), then choose \(K\) from the calibrated table.\\
      Fall back & Use a fixed \(K\) when calibration is insufficient.\\
      \bottomrule
    \end{tabular}
  \end{center}
  \vspace{0.2cm}
  \begin{center}
    Include the controller's own time in every comparison.
  \end{center}
  \slidecite{Untested hypothesis H1 and proposed mechanism M2 in source-notes.md. No gain is assumed.}
\end{frame}

\begin{frame}{A crossed experiment tests load against draft length}
  \vspace{0.05cm}
  \begin{center}
    Proposed grid. Every cell is a planned measurement.
    \vspace{0.25cm}

    \renewcommand{\arraystretch}{1.35}
    \begin{tabular}{rcccc}
      \toprule
      Active requests \(B\) & \(K=1\) & \(K=2\) & \(K=4\) & \(K=8\)\\
      \midrule
      1 & \(\circ\) & \(\circ\) & \(\circ\) & \(\circ\)\\
      4 & \(\circ\) & \(\circ\) & \(\circ\) & \(\circ\)\\
      16 & \(\circ\) & \(\circ\) & \(\circ\) & \(\circ\)\\
      \bottomrule
    \end{tabular}
  \end{center}
  \vspace{0.3cm}
  \begin{itemize}
    \item Calibrate both the adaptive policy and a strong fixed-\(K\) baseline.
    \item Lock both policies, then compare matched held-out load traces.
  \end{itemize}
  \slidecite{Proposed validation V1 in source-notes.md. Grid values are authored test choices, not benchmark data.}
\end{frame}

\begin{frame}{Latency and throughput need a shared evaluation setting}
  \vspace{0.1cm}
  \begin{center}
    \renewcommand{\arraystretch}{1.5}
    \begin{tabular}{lp{8.5cm}}
      \toprule
      Planned metric & Definition\\
      \midrule
      Latency & End-to-end request time, including queueing: p50 and p95 in ms.\\
      Throughput & Emitted output tokens per second.\\
      Acceptance & Accepted draft tokens / proposed draft tokens.\\
      \bottomrule
    \end{tabular}
  \end{center}
  \vspace{0.35cm}
  \begin{itemize}
    \item Match hardware, prompts and output budgets. Include policy overhead.
    \item Repeat paired load traces and report variability.
  \end{itemize}
  \slidecite{Metrics required by P1; proposed definitions and controls V2 in source-notes.md. No outcomes collected.}
\end{frame}

\begin{frame}{The proposal can fail even when drafting is faster}
  \vspace{0.2cm}
  \begin{itemize}
    \item A token-distribution mismatch rejects the policy.
    \item No robust held-out latency-throughput gain over the calibrated fixed baseline weakens the hypothesis.
    \item Controller overhead can erase any benefit from a different length.
  \end{itemize}
  \vspace{0.5cm}
  \begin{center}
    {\large The next evidence is a controlled comparison.}\\[0.25cm]
    Current status: untested proposal.
  \end{center}
  \slidecite{Proposed falsifiers F1 in source-notes.md. Prespecify the operating point before interpreting trade-offs.}
\end{frame}
\end{document}
